\section{Relation to earlier work and a same-object correction}
\label{sec:prior}

\paragraph{Multiplicative shifts and leading geometry.}
The multiplicative golden-mean system and its leading counts and dimensions
belong to the early literature \citep{fan2012level}.  The broader
multiplicative-integer framework and leading Hausdorff/Minkowski theory were
developed by \citet{kenyon2012hausdorff}.  Direct admissible-chain pattern
products, chain-length densities, entropy, and leading error control are
likewise established in \citet{ban2019pattern}.  Recent affine extensions
broaden that leading dimension theory \citep{ban2025affine}.  We use these
objects and formulas as inputs and assign them zero contribution credit.

\paragraph{Digital fluctuations as a neighboring method family.}
Summatory functions of digit-additive sequences often carry periodic or
fractal fluctuations.  The literature represented by
\citet{madritsch2012summatory} is therefore a methodological neighbor.
Its real periodic fluctuations do not, by themselves, state the exact
valuation increment, the continuous inverse-limit remainder, or the
same-object accumulation theorem proved below.  The distinction matters:
our coordinate \(x\bmod q^v\) is first selected in the finite quotient and
then evaluated as a real number; no fractional-part substitution or
\(q\)-adic-valued analytic function is involved.

\paragraph{Boundary complexity.}
\citet{ban2023boundary} introduce boundary complexity and surface entropy
for multiplicative integer systems.  Their valid leading and boundary
results, together with that terminology, are established prior work.  The present
paper asks a narrower question: an exact bounded remainder at every integer
cutoff and the closure of all its subsequential values.  The ownership
separation is summarized in \cref{tab:ownership}.

\begin{table}[t]
  \centering
  \caption{Direct attribution of ingredients and results.  Previously
  established components receive zero credit here.  ``Proved here'' records
  mathematical location only and is not a priority assertion.}
  \label{tab:ownership}
  \small
  \begin{tabularx}{\textwidth}{>{\raggedright\arraybackslash}X
      >{\raggedright\arraybackslash}p{3.8cm}
      >{\raggedright\arraybackslash}p{3.1cm}}
    \toprule
    Component & Source/status & Credit in this paper \\
    \midrule
    Multiplicative SFT, \(q\)-adic chain partition, chain product
      & standard/elementary & zero \\
    Fibonacci word counts; leading entropy and shift dimensions
      & \citet{fan2012level,kenyon2012hausdorff,ban2019pattern} & zero \\
    Boundary-complexity and surface-entropy framework
      & \citet{ban2023boundary} & zero \\
    Affine multiplicative-shift dimensions
      & \citet{ban2025affine} & zero \\
    Exact valuation increment and continuous real boundary on \(\mathbb Z_q\)
      & proved in \cref{sec:exact} & theorem result; no priority claim \\
    Golden full-tail separation, image dimension, radial coefficient
      & proved in \cref{sec:golden,sec:natural} & theorem result; golden case only \\
    \bottomrule
  \end{tabularx}
\end{table}

\subsection{The author-manuscript specialization}
\label{sec:bhl-correction}

The checked artifact is the author manuscript
\href{https://arxiv.org/abs/2210.09115v1}{arXiv:2210.09115v1}, submitted
17 October 2022.  Theorem 3.3(2), the proof to which it reduces, and
Remark 3.4 in that artifact give the relevant locator and hypotheses.  The
proof of Theorem 3.3 says its proofs are similar to that of Theorem 3.1 after
the corresponding replacements; the proof of Theorem 3.1(2) explicitly specializes
\(\Omega\) to a one-dimensional mixing SFT with transition matrix \(A\) and
Perron eigenvalue \(\lambda_A\).  Remark 3.4 then takes spatial dimension one,
\(p\ge2\), \(k\in\mathbb N\), and \(N=p^{kn}\), and displays
\begin{equation}
  \log\lvert\mathcal P([1,p^{kn}],X_{\Sigma_A}^{p})\rvert
  -p^{kn}h
  =-(1-p^{-1})\log(\lambda_A)\,kn+o(kn).
  \label{eq:bhl-displayed}
\end{equation}
The dictionary needed to compare \eqref{eq:bhl-displayed} with our object is
explicit:
\begin{center}
\small
\begin{tabular}{llll}
  \toprule
  author-manuscript symbol & present symbol & specialization & meaning \\
  \midrule
  \(X_{\Sigma_A}^{p}\) & \(X_A^{(q)}\) & \(p=q\) & multiplicative SFT \\
  \(\lvert\mathcal P([1,p^{kn}],X_{\Sigma_A}^{p})\rvert\)
    & \(Z(N)\) & \(N=p^{kn}\) & prefix count \\
  \(\lambda_A\) & \(\rho(A)\) & same adjacency & Perron eigenvalue \\
  \(h\) & \(h\) & same normalization & prefix entropy \\
  \bottomrule
\end{tabular}
\end{center}

Take the full shift \(A=J_D\), the \(D\times D\) all-one matrix, with
\(D>1\).  This matrix is positive, hence primitive, and its edge shift is
mixing; it therefore satisfies the proof specialization just stated.  Every
site is free, hence
\[
  W_\ell=D^\ell,\qquad Z(N)=D^N,\qquad h=\log D.
\]
The left side of \eqref{eq:bhl-displayed} is therefore identically zero,
whereas its displayed main term equals
\(-(1-p^{-1})\log D\,kn\), which is nonzero.  More generally, the exact
identity in \cref{thm:qadic-boundary} gives \(d_v=0\) and \(E(N)=0\) for
this control.  Thus the displayed one-dimensional specialization cannot
hold with the stated quantifiers.  The bounded order-one identity below is
the compatible same-object correction under our present primitive
one-dimensional definition.

\begin{quote}\small
We checked only the author manuscript arXiv:2210.09115v1, submitted
17 October 2022.  We did not line-check the version of record or any
erratum, and make no claim about their displayed formulas.
\end{quote}
Accordingly, we do not transfer \eqref{eq:bhl-displayed} to the journal
version.  We preserve every valid leading and boundary result of
\citet{ban2023boundary}, assign this correction zero prior-ownership credit,
and make no first-discovery claim.  A primary source containing the exact
\(\mathbb Z_q\) extension and complete accumulation theorem, or the same
golden separation theorem, would trigger the conditional duplicate stop
recorded in \cref{app:route-chronology}.
